\chapter{Ecuații diferențiale de ordin superior}

\index{ecuații!de ordin superior}

\section{Ecuații liniare de ordinul $n$ cu coeficienți constanți}

Forma generală este:
\[
  L[y] = a_0 y^{(n)} + a_1 y^{(n-1)} + \dots + a_n y = f(x),
\]
cu $a_i \in \RR$ constante.

Dacă avem $f(x) = 0$, atunci ecuația se numește \emph{omogenă}.

Avem o metodă algebrică de a rezolva această ecuație, folosind:

\begin{definition}\label{def:pol-car-ec}
  Se numește \emph{polinomul caracteristic} atașat ecuației omogene $L[y] = 0$
  polinomul:
  \[
    F(r) = a_0 r^n + a_1 r^{n-1} + \dots + a_n,
  \]
  iar $F(r) = 0$ se numește \emph{ecuația caracteristică} atașată ecuației diferențiale.
\end{definition}

Folosind această noțiune, putem rezolva direct ecuația, ținînd seama de următoarele
cazuri posibile:

\begin{theorem}\label{thm:ec-car-dif}
  \begin{enumerate}[(1)]
  \item Dacă ecuația caracteristică $F(r) = 0$ are rădăcini reale și distincte $r_i$,
    atunci un sistem fundamental de soluții este dat de:
    \[
      \big\{ y_i(x) = e^{r_i x} \big\}_{i=1,\dots,n}.
    \]
  \item Dacă printre rădăcinile lui $F(r)$ există și rădăcini multiple, de exemplu
    $r_1$, cu ordinul de multiplicitate $p$, atunci pentru această rădăcină avem
    $p$ soluții liniar independente (pe lîngă celelalte):
    \[
      y_1(x) = e^{r_1 x}, y_2(x) = xe^{r_1 x}, \dots, y_p(x) = x^{p-1} e^{r_1 x}.
    \]
  \item Dacă printre rădăcinile ecuației caracteristice avem și rădăcini complexe,
    de exemplu $r = a + ib$ și $\overline{r} = a - ib$, atunci fiecărei perechi
    de rădăcini complexe conjugate îi corespund două soluții liniar independente
    (pe lîngă celelalte):
    \[
      y_1(x) = e^{ax} \cos bx, \quad y_2(x) = e^{ax} \sin bx.
    \]
  \item Dacă ecuația caracteristică are rădăcini complexe ca mai sus, cu ordinul
    de multiplicitate $p$, atunci lor le corespund $2p$ soluții liniar independente:
    \[
      \big\{ y_j(x) = x^{j-1} e^{ax} \cos bx \big\}_{j=1, \dots p}, \quad
      \big\{ y_k(x) = x^{k-p-1} e^{ax} \sin bx\big\}_{k=p+1, \dots, 2p}.
    \]
  \end{enumerate}
\end{theorem}

De exemplu: $y^{\prime\prime} - y = 0$, cu condițiile $y(0) = 2$ și $y'(0) = 0$.

\emph{Soluție:} Ecuația caracteristică este $r^2 - 1 = 0$, deci $r_{1,2} = \pm 1$.
Sîntem în primul caz al teoremei, deci un sistem fundamental de soluții este:
\[
  y_1(x) = e^{-x}, \quad y_2(x) = e^x.
\]
Soluția generală este $y(x) = c_1 e^{-x} + c_2 e^x$.

Folosind condițiile Cauchy, obținem $c_1 = c_2 = 1$, deci soluția particulară este:
\[
  y(x) = e^{-x} + e^x.
\]

\vspace{1cm}

\begin{remark}\label{rk:neomogen}
  Dacă ecuația inițială $L[y] = f(x)$ \emph{nu} este omogenă, putem rezolva folosind
  \emph{metoda variației constantelor (Lagrange)}.
\end{remark}

În acest caz, metoda variației constantelor presupune următoarele etape.
Fie ecuația neomogenă scrisă în forma:
\[
  \sum_{k = 0}^n a_k y^{(k)} = f(x).
\]

Pentru simplitate, vom presupune $n = 2$ și fie o soluție particulară de forma:
\[
  y_p(x) = c_1(x) y_1 + c_2(x) y_2.
\]

Atunci, prin metoda variației constantelor, vom determina funcțiile $c_1(x), c_2(x)$
ca soluții ale sistemului algebric:
\[
  \begin{cases}
    c'_1(x) y_1 + c'_2(x) y_2 &= 0 \\
    c'_1(x) y'_1 + c'_2(x) y'_2 &= \dfrac{f(x)}{a_n(x)}
  \end{cases}
\]

\vspace{1cm}

De exemplu: $y^{\prime \prime} + 3y' + 2y = \frac{1}{e^x + 1}$.

\emph{Soluție:} Asociem ecuația omogenă, care are ecuația caracteristică
$r^2 + 3r + 2 = 0$, cu rădăcinile $r_1 = -1, r_2 = -2$. Așadar, soluția generală
a ecuației omogene este:
\[
  \overline{y}(x) = c_1 e^{-x} + c_2 e^{-2x}.
\]

Determinăm o soluție particulară $y_p(x)$ cu ajutorul metodei variației constantelor.
Mai precis, căutăm:
\[
  y_p(x) = c_1(x) e^{-x} + c_2(x) e^{-2x}.
\]

Înlocuind în ecuația neomogenă, rezultă sistemul:
\[
  \begin{cases}
    c_1'(x) e^{-x} + c_2'(x) e^{-2x} &= 0 \\
    -c_1'(x) e^{-x} - 2c_2'(x) e^{-2x} &= \frac{1}{1 + e^x}
  \end{cases}
\]

Rezolvînd ca pe un sistem algebric, obținem:
\[
  c_1'(x) = \frac{e^x}{1 + e^x}, \quad c_2'(x) = -\frac{e^{2x}}{1 + e^x},
\]
de unde obținem:
\[
  c_1(x) = \ln(1 + e^x), \quad c_2(x) = -e^x + \ln(1 + e^x).
\]

În fine, înlocuim în soluția particulară $y_p$ și apoi în cea generală,
$y(x) = \overline{y}(x) + y_p(x)$.

\vspace{1cm}

Un alt exemplu: $y^{\prime\prime} - y = 4e^x$.

\emph{Soluție:} Ecuația caracteristică $r^2 - 1$ are rădăcinile $r_{1,2} = \pm 1$,
deci soluția generală a ecuației liniare omogene este $y_0(x) = c_1 e^x + c_2e^{-x}$.

Căutăm acum o soluție particulară a ecuației neomogene, folosind metoda variației
constantelor. Așadar, căutăm $y_p(x) = c_1(x) e^x + c_2(x) e^{-x}$. Din condiția
ca $y_p$ să verifice ecuația liniară neomogenă, obținem sistemul:
\[
  \begin{cases}
    c_1'(x) e^x + c_2'(x) e^{-x} = 0 \\
    c_1'(x) e^x - c_2'(x) e^{-x} = 4e^x
  \end{cases}
\]
Rezultă $c'_1(x)e^x = 2e^x$, deci $c'_1(x) = 2 \Rightarrow c_1(x) = 2x$, iar
$c'_2(x) e^{-x} = -2e^x$, deci $c_2(x) = -e^{2x}$.

În fine, soluția particulară este $y_p(x) = 2xe^x - e^x$, iar soluția generală:
\[
  y(x) = y_0(x) + y_p(x) = c_1e^x + c_2 e^{-x} + e^x(2x - 1).
\]

%%%%%%%%%%%%%%%%%%%%%%%%%%%%%%%%%%%%%%%%%%%%%%%%%%%%%%%%%%%%%%%%%%%%%% 

\section{Exerciții}

1. Rezolvați următoarele ecuații diferențiale liniare de ordin superior:
\begin{enumerate}[(a)]
\item $ y^{(3)} + 4y^{(2)} + 3y^{(1)} = 0 $;
\item $ y^{(2)} + 4y^{(1)} + 4y = 0 $;
\item $ y^{(3)} + y = 0 $;
\item $ y^{(4)} + 4y^{(2)} = 0 $;
\item $ y^{(2)} - 2y^{(1)} + y = \dfrac{1}{x} e^x $;
\item $ y^{(2)} + y = \dfrac{1}{\cos x} $;
\item $ y^{(3)} + y^{(1)} = \tan x $;
\item $ y^{(3)} + 2y^{(2)} = x + 2 $.
\end{enumerate}

%%%%%%%%%%%%%%%%%%%%%%%%%%%%%%%%%%%%%%%%%%%%%%%%%%%%%%%%%%%%%%%%%%%%%% 

{\color{red}
  \begin{remark}\label{rk:tipuri}
    Tipurile de ecuații pe care le regăsiți mai jos sînt clasificate pentru
    că au forme generale comune. Dar nu este necesar să învățați numele
    acestor tipuri. Concentrați-vă pe metodele de rezolvare și veți constata
    că, indiferent de tipul căreia aparține ecuația, metoda de rezolvare
    este intuitivă. Așadar, nu mai este neapărat cazul obligatoriu, ca pentru
    ecuațiile de ordinul întîi, ca la tipul X de ecuație să se aplice exact
    metoda corespunzătoare.
  \end{remark}
}


\section{Ecuații rezolvabile prin cuadraturi}

\indent\indent Acestea sînt ecuații care se pot rezolva prin integrări
succesive. Astfel, forma lor generală este $ y^{(n)} = f(x) $, în varianta
neomogenă. Soluția generală va depinde de $ n $ constante arbitrare,
rezultate în urma integrărilor succesive.

\textbf{Exemplu:}
\[
  y^{\prime\prime} = \arcsin x + \frac{x}{\sqrt{1 - x^2}} - \frac{\pi}{6}.
\]
\emph{Soluție:} Integrăm succesiv și obținem:
\begin{align*}
  y^{\prime} &= x \arcsin x - \dfrac{\pi}{6} x + c_1 \\
  y &= \dfrac{x^2}{2} \arcsin x + \dfrac{1}{4}x \sqrt{1 - x^2} - %
      \dfrac{1}{4} \arcsin x - \dfrac{\pi}{12} x^2 + c_1 x + c_2.
\end{align*}

Dacă nu se dau condiții inițiale, soluția rămîne exprimată ca mai sus,
adică în funcție de constantele $ c_1 $ și $ c_2 $.

\vspace{1cm}

Cazul omogen se rezolvă și mai simplu: dacă $ y^{(n)} = 0 $, atunci
$ y $ este un polinom de $ x $ de gradul $ n - 1 $.

%%%%%%%%%%%%%%%%%%%%%%%%%%%%%%%%%%%%%%%%%%%%%%%%%%%%%%%%%%%%%%%%%%%%%% 

\section{Ecuații de forma $ f(x, y^{(n)}) = 0 $}

Avem două cazuri care trebuie tratate distinct:
\begin{enumerate}[(a)]
\item Dacă ecuația poate fi rezolvată în raport cu $ y^{(n)} $, adică dacă
  $ \dfrac{\partial f}{\partial y^{(n)}} \neq 0 $, atunci obținem una
  sau mai multe ecuații ca în secțiunea anterioară;
\item Dacă ecuația nu este rezolvabilă cu ajutorul funcțiilor elementare
  în raport cu $ y^{(n)} $, dar cunoaștem o reprezentare parametrică
  a curbei $ f(u, v) = 0 $, adică $ u = g(t), v = h(t) $, cu
  $ t \in [\alpha, \beta] $, atunci soluția generală se poate da sub formă
  parametrică:
  \[
    \begin{cases}
      x &= g(t) \\
      dy^{(n-1)} &= y^{(n)} dx = h(t) g^{\prime}(t) dt
    \end{cases},
  \]
  de unde putem obține succesiv:
  \[
    \begin{cases}
      y^{(n-1)}  &= h_1(t, c_1) \\
      \dotfill \\
      y(t) &= h_n(t, c_1, \dots, c_n).
    \end{cases}
  \]
\end{enumerate}

\textbf{Exemplu:}
\[
  x - e^{y^{\prime\prime}} + y^{\prime\prime} = 0.
\]
\emph{Soluție:} Putem nota $ y^{\prime\prime} = t $ și atunci $ x(t) = e^t - t $.
Deoarece avem $ dy^{\prime} = y^{\prime\prime} dx $, rezultă:
\[
  dy^{\prime} = t(e^t - 1) dt \Rightarrow y^{\prime} = -\dfrac{t^2}{2} + te^t - e^t + c_1.
\]
Mai departe, $ dy = y^{\prime} dx $, deci:
\[
  dy = \Big( - \dfrac{t^2}{2} + te^t - e^t + c_1 \Big)(e^t - 1) dt.
\]
În fine, soluția este:
\[
  y(t) = e^{2t} \Big( \dfrac{t}{2} - \dfrac{3}{4} \Big) + e^t \Big( - \dfrac{t^2}{2} + %
  1 + c_1 \Big) + \dfrac{t^3}{6} - c_1 t + c_2.
\]

%%%%%%%%%%%%%%%%%%%%%%%%%%%%%%%%%%%%%%%%%%%%%%%%%%%%%%%%%%%%%%%%%%%%%% 

\section{Ecuații de forma $ f(y^{(n-1)}, y^{(n)}) = 0 $}

Din nou, distingem două cazuri:
\begin{enumerate}[(a)]
\item Dacă ecuația este rezolvabilă prin funcții elementare în raport cu $ y^{(n)} $,
  atunci putem nota $ z = y^{(n-1)} $ și avem $ z^{\prime} = f(z) $. Ecuația devine cu
  variabile separabile pentru $ z $ și se rezolvă corespunzător, conducînd la
  $ z = y^{(n-1)} = f_1(x, c_1) $, care este de tipul anterior;
\item Dacă ecuația nu este rezolvabilă prin funcții elementare în raport cu
  $ y^{(n)} $, dar cunoaștem o reprezentare parametrică de forma:
  \[
    y^{(n-1)} = h(t), \quad y^{(n)} = g(t), \quad t \in [\alpha, \beta],
  \]
  atunci putem folosi $ dy^{(n-1)} = y^{(n)} dx $, pentru a obține soluția
  pas cu pas, sub formă parametrică:
  \begin{align*}
    x &= \ds\int\frac{h^{\prime}}{g} dt + c_1 \\
    y^{(n-1)} &= h(t) \\
    dy^{(n-2)} &= h(t) = \dfrac{hh^{\prime}}{g} dt \\
    y^{(n-2)} &= \ds\int \frac{hh^{\prime}}{g} dt + c_2 \\
      &\vdots \\
    y &= \int y^{\prime} dx + c_n.
  \end{align*}

\end{enumerate}

\textbf{Exemplu:} $ y^{\prime\prime} = - \sqrt{1 - y^{^{\prime}2}} $.
\emph{Soluție:} Facem schimbarea de funcție $ z(x) = y^{\prime} $ și ecuația devine:
\[
  - \dfrac{dz}{\sqrt{1 - z^2}} = dx, \quad |z| < 1.
\]
Soluția generală este: $ \arccos z = x + c_1 $, deci $ z = \cos(x + c_1) $.

Mai departe, obținem $ y(x) = \sin(x + c_1) + c_2 $.

De asemenea, trebuie să mai remarcăm și soluțiile particulare $ y = \pm x + c $.

%%%%%%%%%%%%%%%%%%%%%%%%%%%%%%%%%%%%%%%%%%%%%%%%%%%%%%%%%%%% 

\section{Ecuații de forma $ f(x, y^{(k)}, \dots, y^{(n)}) = 0 $}

Ecuația se rezolvă cu schimbarea de funcție $ y^{(k)} = z(x) $ și rezultă
o ecuație de ordin $ n - k $:
\[
  f(x, z, z^{\prime}, \dots, z^{(n-k)}) = 0,
\]
pe care o integrăm.

\textbf{Exemplu:} $ (1 + x^2)y^{\prime\prime} = 2xy^{\prime} $.

\emph{Soluție:} Cu schimbarea de funcție $ y^{\prime} = z(x) $, obținem ecuația:
\[
  \dfrac{z}{z^{\prime}} = \dfrac{2x}{1 + x^2},
\]
iar apoi, prin integrare, rezultă $ z = y^{\prime} = c_1(1 + x^2) $. În fine:
\[
  y(x) = c_1 x + c_1 \dfrac{x^3}{3} + c_2.
\]
\vspace{1cm}

\textbf{Exemplu 2:} $ y \cdot y^{\prime\prime} - yy^{^{\prime}2} = 0 $.

\emph{Soluție:} Notăm $ y^{\prime} = z $ și obținem:
\[
  y\cdot z^{\prime} - yz^2 = 0 \Rightarrow y(z^{\prime} - z^2) = 0,
\]
de unde rezultă $ y = 0 $ sau $ z^{\prime} - z^2 = 0 $.

Din a doua variantă, deducem succesiv:
\[
  \dfrac{dz}{dx} = z^2 \Rightarrow \dfrac{dz}{z^2} = dx \Rightarrow %
  -\dfrac{1}{z} = x + c_1.
\]
Mai departe:
\[
  -\dfrac{1}{y^{\prime}} = x + c_1 \Rightarrow y^{\prime} = \dfrac{-1}{x + c_1} %
  \Rightarrow y = -\ln(x + c_1) + c_2.
\]


%%%%%%%%%%%%%%%%%%%%%%%%%%%%%%%%%%%%%%%%%%%%%%%%%%%%%%%%%%%%%%%%%%%%%% 

\section{Ecuații de forma $ f(y^{\prime}, y^{\prime\prime}, \dots, y^{(n)}) = 0 $}

Să remarcăm că astfel de ecuații nu conțin pe $ y $. Deci putem lua pe $ y^{\prime} $
ca variabilă independentă și pe $ y^{\prime\prime} $ ca fiind funcție de $ y^{\prime} $.
Astfel, reducem discuția la un caz anterior.

\textbf{Exemplu:} $ x^2 y^{\prime\prime} = y^{^{\prime}2} - 2xy^{\prime} + 2x^2 $.

\emph{Soluție:} Facem schimbarea de funcție $ y^{\prime} = z(x) $ și obținem o
ecuație Riccati:
\[
  z^{\prime} = \dfrac{z^2}{x^2} - 2 \cdot \dfrac{z}{x} + 2.
\]
Observăm soluția particulară $ z_p = x $ și integrăm, cu $ z = x + \dfrac{1}{u(x)} $.

Obținem succesiv:
\begin{align*}
  z(x) &= x + \dfrac{c_1x}{x + c_1} \Rightarrow \\
  y^{\prime}(x) &= x + \dfrac{c_1x}{x + c_1} \Rightarrow \\
  y(x) &= \dfrac{x^2}{2} + c_1 x - c_1^2 \ln | x + c_1 | + c_2,
\end{align*}

%%%%%%%%%%%%%%%%%%%%%%%%%%%%%%%%%%%%%%%%%%%%%%%%%%%%%%%%%%%%%%%%%%%%%% 

\section{Ecuații autonome (ce nu conțin pe $ x $)}

În cazul acestor ecuații, putem micșora ordinul cu o unitate, notînd
$ y^{\prime} = p $ și luăm pe $ y $ variabilă independentă.

\begin{remark}\label{rk:aut-const}
  Există posibilitatea de a pierde soluții de forma $ y = c $ prin această
  metodă, deci trebuie verificat dacă este cazul separat.
\end{remark}

\textbf{Exemplu:} $ 1 + y^{^{\prime}2} = 2yy^{\prime} $.
\emph{Soluție:} Putem lua $ y^{\prime} = p $ drept funcție, iar pe $ y $ drept
variabilă independentă. Rezultă:
\[
  y^{\prime\prime} = \dfrac{d}{dx} \Big( \dfrac{dy}{dx} \Big) = p \dfrac{dp}{dy}.
\]
Atunci ecuația devine:
\[
  \dfrac{2p dp}{1 + p^2} = \dfrac{dy}{y} \Rightarrow %
  y = c_1(1 + p^2).
\]
Acum trebuie să obținem pe $ x $ ca funcție de $ p $ și $ c_1 $.
Deoarece $ dx = \dfrac{1}{p} dy $, iar $ dy = 2c_1 pdp $, rezultă
$ dx = 2c_1 dp $. Deci $ x = 2c_1p + c_2 $, iar soluția generală devine
$ x(p) = 2c_1 p + c_2 $. Din expresia lui $ y $ de mai sus, rezultă:
\[
  y = c_1 + \dfrac{(x - c_2)^2}{4c_1}.
\]

\vspace{1cm}

\textbf{Exemplu 2:} $ y^{\prime\prime} = y^{^{\prime}2} = 2e^{-y} $.

\emph{Soluție:} Notăm $ y^{\prime} = p $ și luăm ca necunoscută $ p = p(y) $.
Rezultă:
\[
  y^{\prime\prime} = \dfrac{dy^{\prime}}{dx} = p^{\prime}p \Rightarrow p^{\prime}p + p^2 = 2e^{-y}.
\]
Dacă notăm $ p^2 = z $, obținem o ecuație liniară neomogenă:
\[
  z^{\prime} + 2z = 4e^{-y},
\]
ce are ca soluție generală $ z(y) = c_1 e^{-2y} + 4e^{-y} $.

Revenim la $ y $ și avem:
\[
  z = p^2 = y^{^{\prime}2} \Rightarrow y^{\prime} = \pm \sqrt{c_1 e^{-2y} + 4e^{-y}},
\]
adică:
\[
  \dfrac{dy}{\pm \sqrt{c_1 e^{-2y} + 4e^{-y}}} = dx \Rightarrow %
  x + c_2 = \pm \dfrac{1}{2} \sqrt{c_1 + 4e^y}.
\]
Echivalent, putem scrie soluția și în forma implicită:
\[
  e^y + \dfrac{c_1}{4} = (x + c_2)^2.
\]

%%%%%%%%%%%%%%%%%%%%%%%%%%%%%%%%%%%%%%%%%%%%%%%%%%%%%%%%%%%%%%%%%%%%%% 

\section{Ecuații Euler}
\index{ecuații!Euler}
O ecuație Euler are forma generală:
\[
  \sum_{k = 0}^n a_k x^k y^{(k)} = f(x),
\]
pentru varianta omogenă și $ f(x) = 0 $ pentru varianta neomogenă, cu $ a_i $
constante reale.

Ecuațiile Euler se pot transforma în ecuații cu coeficienți constanți prin notația
(schimbarea de variabilă) $ |x| = e^t $.

De remarcat faptul că, deoarece funcția necunoscută este $ y = y(x) $, pentru
a obține $ y^{\prime} = \dfrac{dy}{dx} $ trebuie să aplicăm regula de derivare
a funcțiilor compuse. Folosind notația de la fizică $ \dot{y} = \dfrac{dy}{dt} $,
putem scrie:
\[
  \dfrac{dy}{dx} = \dfrac{dy}{dt} \cdot \dfrac{dt}{dx} = \dot{y} \cdot e^{-t}.
\]
Mai departe, de exemplu, pentru derivatele superioare:
\begin{align*}
  \dfrac{d^2 y}{dx^2} &= \dfrac{d}{dt} \Big( \dot{y} \cdot e^{-t} \Big) %
                        \cdot \dfrac{dt}{dx} \\
                      &= e^{-2t} (\ddot{y} - \dot{y}).
\end{align*}

\vspace{1cm}

\textbf{Exemplu:} $ x^2 y^{\prime\prime} + xy^{\prime} + x = x $.

\emph{Soluție:} Facem substituția $ |x| = e^t $ și, folosind calculele
de mai sus, obținem:
\[
  e^{2t} \cdot e^{-2t} (\ddot{y} - \dot{y}) - e^t \cdot e^{-t} + y(t) = e^t.
\]
Echivalent: $ \ddot{y} - 2\dot{y} + y = e^t $.

Aceasta este o ecuație liniară de ordinul al doilea, neomogenă. Asociem ecuația
algebrică $ r^2 - 2r + 1 = (r - 1)^2 = 0 $, deci soluția generală a variantei
omogene este:
\[
  y(t) = c_1 e^t + c_2 te^t.
\]

Folosind metoda variației constantelor, obținem succesiv:
\begin{align*}
  y(t) &= c_1 e^t + c_2 t e^t + \dfrac{t^2 e^t}{2} \Rightarrow \\
  y(x) &= c_1 x + c_2 x \ln x + \dfrac{x \ln^2 x}{2}.
\end{align*}


%%%%%%%%%%%%%%%%%%%%%%%%%%%%%%%%%%%%%%%%%%%%%%%%%%%%%%%%%%%%%%%%%%%%%% 

\section{Exerciții}

1. Rezolvați următoarele ecuații Euler:
\begin{enumerate}[(a)]
\item $ x^2 y^{(2)} - 3xy^{\prime} + 4y = 5x, \quad x > 0 $;
\item $ x^2 y^{(2)} - xy^{\prime} + y = x + \ln x, \quad x > 0 $;
\item $ 4(x + 1)^2 y^{(2)} + y = 0, x > -1 $;
\item $ xy^{\prime\prime} + 3y^{\prime} = 1 $;
\item $ x^3 y^{\prime\prime\prime} + x^2y^{\prime\prime} - 2xy^{\prime} + 2y = \dfrac{1}{x^2} $.
\end{enumerate}


\vspace{1cm}

2. Rezolvați următoarele ecuații:
\begin{enumerate}[(a)]
\item $ y^{(3)} = \sin x + \cos x $;
\item $ y^{(3)} + y^{(2)} = \sin x $;
\item $ xy^{(2)} + (x - 2)y^{(1)} - 2y = 0 $;
\item $ x^2 y^{(2)} = y^{\prime 2} $;
\item $ y^{(4)} - 3y^{(3)} + 3y^{(2)} - y^{(1)} = 0 $;
\item $ y^{(2)} - y^{(1)} - 2y = 3e^{2x} $;
\item $ 4(x + 1)^2 y^{(2)} + y = 0 $;
\end{enumerate}

Amintim, pentru ecuații liniare și neomogene de ordinul întîi, cu forma generală:
\[
  y' = P(x) y + Q(x),
\]
avem soluțiile:
\begin{align*}
  y_p &= c \cdot \exp\Big( - \int_{x_0}^x P(t) dt \Big) \\
  y_g &= \Big( c + \int Q(t) \cdot \exp \Big( \int_{x_0}^t P(s) ds \Big) dt \Big) \cdot \exp \Big(- \int_{x_0}^x P(t) dt \Big).
\end{align*}

Puteți folosi formulele de mai sus, sau, \emph{preferabil}, puteți rezolva direct,
ca în seminarul anterior.

\emph{Indicații:}
\begin{enumerate}[(a)]
\item Integrări succesive;
\item Notăm $ y^{(2)} = z $, deci ecuația devine $ z^\prime + z = \sin x $, care este
  o ecuație liniară și neomogenă, de ordinul întîi. Soluția generală este:
  \[
    z(x) = e^{-x}\Big( c + \dfrac{e^x}{2} (\sin x + \cos x) \Big),
  \]
  iar apoi $ y^\prime(x) = \ds\int z(x) dx $ și $ y(x) = \ds\int y^\prime (x) dx $ etc.
\item Notăm $ z = y^\prime + y $, deci ecuația devine $ xz^\prime - 2z = 0 $, care
  are soluția evidentă $ z = c_1 x^2 $, din care obtinem $ y^\prime + y = c_1 x^2 $, care
  este o ecuație liniară și neomogenă, cu soluția generală:
  \[
    y = c_2 e^{-x} + c_1 x^2 - 2c_1 x + 2c_1.
  \]
\item Notăm $ y^\prime = z $ și obținem $ x^2 z^\prime = z^2 $. Observăm soluția particulară
  $ z = 0 $, deci $ y = c $, iar în celelalte cazuri, avem:
  \[
    x^2 \dfrac{dz}{dx} = z^2 \Rightarrow \dfrac{dx}{x^2} = \dfrac{dz}{z^2}.
  \]
  Aceasta conduce la $ \dfrac{1}{z} = \dfrac{1}{x} + c $ și ne întoarcem
  la $ y $:
  \[
    dx = dy \Big( \dfrac{1}{x} + c \Big) \Rightarrow \dfrac{xdx}{cx + 1} = dy,
  \]
  care, prin integrare, conduce la:
  \[
    y = \dfrac{x}{c} - \dfrac{1}{c^2} \ln(cx + 1) + c_1,
  \]
  pentru $ c \neq 0 $ și $ y^\prime = x $ dacă $ c = 0 $, adică $ y = \dfrac{x^2}{2} + c_1 $.
\item Ecuația caracteristică este:
  \[
    r^4 - 3r^3 + 3r^2 - r = 0 \Rightarrow r(r^3 - 3r^2 + 3r - 1) = 0,
  \]
  de unde observăm soluțiile $ r = 0, r = 1 $ etc.
\item Ecuația caracteristică este $ r^2 - r - 2 = (r + 1)(r - 2) = 0 $.
\item Este o ecuație Euler în raport cu $ t = x + 1 $.
\end{enumerate}


%%% Local Variables:
%%% mode: latex
%%% TeX-master: "../edtc-18-19"
%%% End:
